\documentclass[10pt,a4paper]{amsart}
\usepackage[margin=19mm]{geometry}
\usepackage{amsmath,amssymb,mathtools}
\usepackage[scr=rsfs,cal=euler]{mathalfa}
\usepackage{xfrac}
\usepackage[unicode,hidelinks,bookmarks=false]{hyperref}
\newcommand{\ie}{i.e.\ }
\newcommand{\cf}{cf.\ }
\newcommand{\resp}{respectively\ }
\newcommand{\Subsection}{\subsection}
\providecommand{\nobreakdash}{\nobreak\hbox{-}}
\setlength{\emergencystretch}{2em}
\multlinegap0pt
\usepackage{amssymb}
\usepackage{dsfont}
\usepackage{enumerate}
\usepackage{mathtools}
\usepackage{xcolor}
\newtheorem{proposition}{Proposition}
\title{On Conjecture of Binomial Edge Ideals of Linear Type}
\author{Marie Amalore Nambi, Neeraj Kumar}
\date{Source: arXiv:2406.05960v4 (CC BY 4.0)}
\begin{document}
\maketitle

\noindent\emph{Source and licence.} On Conjecture of Binomial Edge Ideals of Linear Type. Version arXiv:2406.05960v4 (CC BY 4.0), distributed by arXiv under the Creative Commons Attribution 4.0 International licence, http://creativecommons.org/licenses/by/4.0/, as recorded in the arXiv licence metadata for this exact version and shown on the abstract page https://arxiv.org/abs/2406.05960v4. Full licence notice: \texttt{licenses/arxiv-2406.05960-CC-BY-4.0.txt}.\par\smallskip

\noindent\textit{Editorial scope.} Counted unit: the article's proposition prop.p-mon (source0.tex lines 271--284). The article's complete original proof of it is delivered with it (source0.tex lines 286--301, one \texttt{proof} environment of 16 lines). No mathematics is written, repaired or re-derived: the statement and the proof are the article's own lines, in the article's own notation, and no retained line is substituted or re-flowed.  No further source-local context is needed: every cross-reference of every retained line resolves inside the two delivered spans.  Nothing was omitted from any retained span, so the package carries no omission record.  No retained line contains a citation, so the wrapper carries no bibliography and no bibliography entry.  No retained line contains a figure environment, an image-inclusion command or drawing code, so no image is delivered and the package declares no asset dependency.  Editorial formatting only, in addition: an AMS \texttt{amsart} preamble that restates the article's own declarations and macros used by the retained lines, namely the environments \texttt{proposition} on separate counters, together with the standard packages those lines need.  The retained environments print in this document as Proposition 1 in order of appearance, whereas the article prints the same items with the numbering of its own full document.  The complete native source of the article as distributed in the arXiv e-print for this exact version is bundled unchanged as \texttt{sources/160526/source0.tex}.
\begin{proposition}\label{prop.p-mon}
Let \(z_1,\ldots,z_n\) be an ordered sequence of monomials in \(\mathbb{K}[x_1,\ldots,x_n]\). Suppose that \(z_i \nmid z_j\) for all \(i \neq j\), and that the following conditions hold:
\begin{equation}\label{cond1}
\gcd(z_i,z_{i_2}) \mid z_{i_1},
\qquad
1 \leq i < i_1 \leq i_2 \leq n,
\end{equation}
\begin{equation}\label{cond2}
\gcd(z_i,z_j^2) = \gcd(z_i,z_j),
\qquad
1 \leq i < j \leq n.
\end{equation}
Then \(z_1,\ldots,z_n\) is a \(p\)-sequence.
\end{proposition}

\begin{proof}
Since $z_i\nmid z_j$ for all $i\neq j$, the sequence $z_1,\ldots,z_n$ is a minimal system of monomial generators. To prove that \(z_1,\ldots,z_n\) is a \(p\)-sequence, it suffices to show that
\[
\gcd(z_i, z_{i_1}z_{i_2}) = \gcd(z_i, z_{i_1})
\quad \text{for all } 1 \leq i < i_1 \leq i_2 \leq n.
\]

Fix \(i < i_1 \leq i_2\). By \eqref{cond1}, \(\gcd(z_i, z_{i_2}) \mid z_{i_1}\); hence, it follows that \(\gcd(z_i, z_{i_1}z_{i_2}) \mid \gcd(z_i, z_{i_1}^2)\). Using \eqref{cond2}, we obtain $\gcd(z_i, z_{i_1}^2) = \gcd(z_i, z_{i_1})$,
and hence
$\gcd(z_i, z_{i_1}z_{i_2}) \mid \gcd(z_i, z_{i_1})$.
On the other hand, it is clear that
$\gcd(z_i, z_{i_1}) \mid \gcd(z_i, z_{i_1}z_{i_2})$.
Therefore,
$\gcd(z_i, z_{i_1}z_{i_2}) = \gcd(z_i, z_{i_1})$,
for all \(1 \leq i < i_1 \leq i_2 \leq n\). This proves that \(z_1,\ldots,z_n\) is a \(p\)-sequence.
\end{proof}

\end{document}
